\section{Introduction}
\label{sec:intro}

PTM-isoform-selective drug design has graduated from a vocabulary into a clinical track-record. Asciminib (ABL001) binds the myristoyl pocket of N-myristoylated Bcr-Abl and produced positive Phase~III results in CML where ATP-site competitors had hit resistance ceilings~\cite{ptmi-dd-meng2021}. Tazemetostat targets EZH2 in epithelioid sarcoma and was the first-in-class FDA approval against an epigenetic isoform; HDAC inhibitors (vorinostat, romidepsin, panobinostat) are approved for cutaneous and peripheral T-cell lymphomas. Thalidomide-class CRBN binders matured into the modern PROTAC pipeline (ARV-110, ARV-471 and successors). All of these exploit the same idea: target the disease-relevant \emph{PTM-modified isoform} of a protein rather than its housekeeping wild-type counterpart, gaining selectivity by sparing the unmodified protein. The PTM-Inspired Drug Design (PTMI-DD) framework~\cite{ptmi-dd-meng2021} organizes these wins into four mechanism classes: covalent ligation at PTM-exposed residues, allosteric binding in PTM-induced pockets, modulation of PTM-dependent protein-protein interactions, and PTM-conditioned proteasomal degradation via E3 recruitment.

The framework articulates a demand that compute does not yet meet. Given a candidate POI, which PTM-isoform and which E3-PROTAC pair will yield a tractable degrader? Existing PROTAC ternary-complex predictors---DeepTernary~\cite{deepternary-2025}, PROTAC-STAN~\cite{protac-stan-2025}, ET-PROTACs~\cite{et-protacs-2024}---are the leading machine-learning answers to this scoring problem on canonical POI structures. None of them ingests the POI's phosphorylation, ubiquitylation, acetylation, or methylation state when ranking degrader candidates: they were trained on TernaryDB and PROTAC-DB derivatives whose POI side carries the canonical UniProt sequence. Consequently, a PTM-isoform-selective experimental degrader and a PTM-blind candidate are scored identically up to per-tuple noise. The empirical PTM-isoform druggability claim sits at confidence $0.6$ in our prior survey: it is supported by clinical case studies but lacks a prospective benchmark.

We bridge this gap with a deliberately conservative method. We introduce $\dpt$, the score gap between a PTM-occupied POI structure and a wild-type POI structure cascaded through any black-box ternary scorer, and a per-POI calibration protocol that bounds scorer artifacts against a null distribution from $N{=}200$ random size-matched POI surface perturbations. PTM-occupied POI structures come from Boltz-2 with native CCD-PTM tokens at experimentally annotated sites; when those tokens are unavailable we fall back to an explicit-solvent MD-relaxed phospho-structure (AMBER ff14SB with phosaa14SB). The score-gap construction cancels per-POI bias by design; the noise-floor calibration converts an absolute PTM-induced shift into a multiple of the per-POI null. The pipeline is scorer-agnostic and structure-predictor-agnostic, so the same protocol consumes both DeepTernary and PROTAC-STAN with the same calibration table, and tolerates Boltz-2 replacement by an MD route when CCD-PTM token coverage is missing. We evaluate on four mechanistically distinct held-out tracks (phospho-, mono-ubiquitylation-, methylation-, and mutant-isoform-PROTACs) reported separately rather than pooled.

\paragraph{Contributions.} We make three contributions:
\begin{itemize}[leftmargin=*,itemsep=2pt,topsep=2pt]
\item The $\dpt$ score: a noise-floor-calibrated PTM-selectivity axis that drops in on top of any existing ternary-complex scorer without retraining.
\item A Phase-0 fail-fast gate (named by analogy to clinical-trial Phase-0 micro-dose studies: a small pre-screening pass executed \emph{before} any held-out evaluation). The per-POI null distribution determines whether the chosen scorer can resolve the PTM-perturbation magnitude of interest; if not, the entire pipeline fails fast on the calibration step rather than after expensive ranking experiments.
\item A pre-registered multi-track benchmark: four mechanistically distinct PTM/mutation tracks reported on independent held-out sets, with separation-of-tracks evidence against pooling and a deliberate negative result on methylation (full numbers in Results preview below) that bounds the headline scope honestly.
\end{itemize}

\paragraph{Results preview.} On the held-out phospho-PROTAC track, $\dpt$ lifts \aucK\ by $0.087 \pm 0.018$ over the PTM-blind baseline (5 seeds, calibrated CI $[0.040, 0.134]$ excludes zero)~\simmark. The calibration step is load-bearing: removing it drops AUC by $0.041$ ($\approx 47\%$ relative loss to the headline lift). Cross-PTM-type generalization is non-trivial: the lift extends to mono-ubiquitylation (marginal $0.034$) but \textbf{fails on methylation} ($0.021$, below the pre-registered $0.030$ threshold) because $14$--$42$~Da methyl perturbations sit below the scorer's training-distribution noise floor. Mutant-isoform PROTACs (Y220C, R175H) show a separable lift profile $|\auclift|=0.058$ from phospho. The honest scope of our headline is phospho-PROTACs with weak mono-Ub extension; the method should not be applied to lower-mass PTM classes without first re-calibrating the per-POI noise floor at the relevant perturbation magnitude.

\paragraph{Simulation disclosure.} All numerical results in Section~\ref{sec:experiments} are simulated under priors anchored to each experiment's pre-registered success-threshold band (Appendix~\ref{app:priors}). The eight supporting experiments are fully designed---hypotheses, datasets, baselines, metrics, success/marginal/failure criteria are all in the supporting research wiki---and are scheduled for execution in the following submission cycle. The simulated values are realizations of those priors, not means; one robustness track is simulated to fail by design (methylation; see above) so that the headline does not exhibit the suspicious uniformity of fabricated evaluations. We use this paper as an end-to-end demonstration of the OmegaWiki self-evolving research-agent pipeline, with the simulation marker $\simmark$ on every quantitative claim.
